\section{The Architect's Question}\label{the-architects-question}

After this chapter we should be able to sit with a customer or
stakeholder who says ``we want to build something with AI'' and turn
that vague ambition into a bounded, measurable, defensible scope. We
will apply the book's core loop in a conversational setting: ask the
questions that surface the real workload (Ch4), the quality bar (Ch3/5),
the latency and cost constraints, and the operational reality --- then
restate the requirement in terms we can architect against. After this
chapter, ``let's go build an AI thing'' becomes ``here are the bounds,
the SLO, and the open questions'' --- and the customer knows exactly
what they asked for.

\section{Concept}\label{concept}

Most AI projects fail in requirements, not in technology. The root cause
is a \textbf{category error}: the customer and the architect mean
different things by the same words. ``Fast,'' ``good,'' ``smart,'' and
``it should just work'' are not contracts; they are vibes. The
architect's first job is to convert vibes into \textbf{quantified,
testable bounds} using the disciplines of the earlier chapters.

The conversion rests on one idea: \textbf{every architectural decision
traces back to a number the customer can see.} Whether the model is 7B
or 70B, whether we self-host or use an API, whether we cache, quantize,
or disaggregate --- all of it is decided by five numbers we can extract
from a customer conversation:

\begin{enumerate}
\def\labelenumi{\arabic{enumi}.}
\tightlist
\item
  \textbf{The traffic} --- how many users, at what concurrency, with
  what peak-to-average ratio (Ch4/Ch17).
\item
  \textbf{The token profile} --- what inputs look like, how long
  contexts are, the input/output ratio (Ch4/Ch6).
\item
  \textbf{The quality bar} --- what ``good enough'' means on \emph{this}
  task, on \emph{their} data (Ch3/Ch5/Ch14).
\item
  \textbf{The latency/cost budget} --- the SLO and the money they can
  spend (Ch6/Ch16).
\item
  \textbf{The operational reality} --- data privacy, geography,
  available staff, time-to-value (Ch13/Ch21/Ch24).
\end{enumerate}

\section{Mental Model}\label{mental-model}

Think of the customer conversation as \textbf{buying information, not
delivering answers.} We are not in the room to impress; we are in the
room to learn the five numbers. Three mental moves keep the conversation
productive:

\begin{itemize}
\tightlist
\item
  \textbf{The Socratic funnel.} Start broad (``what problem are we
  solving?'') and narrow through a scripted set of questions that each
  commit the customer to a concrete bound. Every ``vague'' answer gets
  followed by a request for magnitude: ``roughly how many users?''
  ``what's the slowest acceptable answer?''
\item
  \textbf{The mirror test.} Restate their words back as a quantified
  requirement and check they still agree. ``So: \textasciitilde2,000
  users, under 2 seconds to first token, on company-internal data ---
  did I get that right?'' The restatement is where scope actually gets
  locked.
\item
  \textbf{The two-window honesty.} Separate what the customer
  \emph{thinks} (goals) from what the \emph{system} can guarantee
  (SLOs). Goals are aspirations; SLOs are contracts. Confusing them
  breeds every fee dispute and scope fight.
\end{itemize}

\section{Worked Example}\label{worked-example}

\subsection{The Requirements Dialogue: Turning ``an internal AI
assistant'' Into
Bounds}\label{the-requirements-dialogue-turning-an-internal-ai-assistant-into-bounds}

\textbf{{[}Worked example dialogue --- illustrative, not reference.{]}}
An internal-platform lead says: ``We want an AI assistant for our sales
teams --- it should answer questions from our product docs and CRM data,
and it should be fast.'' {[}2°{]}

\textbf{Architect:} ``\textasciitilde How many people will use it, and
how often?'' \textbf{Lead:} ``\textasciitilde2,000 on the platform;
maybe 10\% are in it at once at peaks.'' \textbf{Architect:} ``So
\textasciitilde200 concurrent at peak, and realistically a fraction of
that firing requests. Roughly how many questions per second at the
busiest minute?'' \textbf{Lead:} ``If everyone on the sales floor asked
at once, maybe 40 a second.'' \textbf{→ Bound 1: \textasciitilde2,000
users, \textasciitilde10 rps average, \textasciitilde40 rps peak.} {[}1P
§14{]}

\textbf{Architect:} ``What does a question look like? Short or with long
context?'' \textbf{Lead:} ``They paste product specs and contracts, so
the prompt can be lengthy --- a few pages.'' \textbf{→ Bound 2: long
\textasciitilde9,200-token inputs with retrieval,
\textasciitilde300-token outputs.} {[}1P §14{]}

\textbf{Architect:} ``What's `fast'? First-word delay, or time to full
answer?'' \textbf{Lead:} ``Feels slow if we wait more than two seconds
to see something start.'' \textbf{→ Bound 3: TTFT ≤ 2 s (p95),
\textasciitilde25 ms/token thereafter.} {[}1P §14{]}

\textbf{Architect:} ``What does `good answer' mean for your team --- and
can we define a handful of sample questions to test against?''
\textbf{Lead:} ``It should be accurate on our docs and not make up
contract terms. We can give the validation team 50 real questions.''
\textbf{→ Bound 4: a 50-question retrieval-QA quality benchmark on their
actual data.} {[}1P §14{]}

\textbf{Architect:} ``Data sensitivity --- can this touch a public
cloud, and how much can it cost?'' \textbf{Lead:} ``It's internal and
regulated; keep it on-prem if at all possible.'' \textbf{→ Bound 5:
operational constraint forces the local/server tier (Ch13).} {[}1P
§14{]}

\textbf{Architect (mirror):} ``So to lock it: \textasciitilde2,000 sales
users, \textasciitilde10 rps average / 40 peak, prompts with
\textasciitilde9,200-token retrieval context, answers starting under 2
s, quality judged on 50 of your real questions, and fully on-prem
because it's regulated. Have I captured it?'' \textbf{Lead:} ``Yes ---
that's exactly it.''

\textbf{Reading the result.} The vague ask collapsed into the exact
canonical workload this book has architected end-to-end --- a 70B-class
model on an 8×H100 on-prem server meets every bound (Ch2-16). The
conversation did not sell a model; it \emph{exposed the numbers} that
then chose the architecture with almost no further input. {[}2°
DERIVED{]} \emph{(worked example, not reference)}

\textbf{\hyperref[tab:23.1]{Table 23-1}} --- The six questions that turn any vague AI ask
into architectable bounds

\begin{longtable}[]{@{}
  >{\raggedright\arraybackslash}p{0.3333\linewidth}
  >{\raggedright\arraybackslash}p{0.3333\linewidth}
  >{\raggedright\arraybackslash}p{0.3333\linewidth}@{}}
\toprule\noalign{}
\begin{minipage}[b]{\linewidth}\raggedright
Question
\end{minipage} & \begin{minipage}[b]{\linewidth}\raggedright
Surface it reveals
\end{minipage} & \begin{minipage}[b]{\linewidth}\raggedright
Canonical answer (this example)
\end{minipage} \\
\midrule\noalign{}
\endhead
\bottomrule\noalign{}
\endlastfoot
\textasciitilde How many users, how concurrent? & Traffic \& concurrency
(Ch4/17) & \textasciitilde2,000 users, \textasciitilde10/40 rps \\
What do inputs look like? & Token profile (Ch4/6) & \textasciitilde9.2K
in / \textasciitilde300 out \\
What's ``fast''? & SLO (Ch6) & TTFT ≤ 2 s, \textasciitilde25 ms/token \\
What's ``good''? & Quality bar (Ch3/14) & 50-question retrieval-QA
set \\
Data sensitive? Cloud OK? & Tier \& ops (Ch13) & on-prem, regulated \\
How much can it cost? & TCO (Ch16) & within self-host budget \\

\label{tab:23.1}\end{longtable}

\section{Measurement}\label{measurement}

The requirements dialogue has its own quality metrics:

\begin{enumerate}
\def\labelenumi{\arabic{enumi}.}
\tightlist
\item
  \textbf{Bound coverage} --- did the conversation surface all five
  numbers (traffic, tokens, quality, latency/cost, ops)? Missing one is
  a future scope surprise.
\item
  \textbf{Customer buy-in} --- did the mirrored restatement get an
  explicit ``yes''? An unconfirmed restatement is not a scope.
\item
  \textbf{Decision lead time} --- how quickly the five bounds let us
  produce a defensible draft architecture. The whole point is that
  extracting numbers, not more meetings, unlocks the design.
\item
  \textbf{Assumption log} --- every guess we made that the customer did
  not explicitly confirm, tracked as {[}VERIFY{]} items, not silent
  defaults.
\end{enumerate}

\section{Common Mistakes}\label{common-mistakes}

\begin{itemize}
\tightlist
\item
  \textbf{Selling before scoping.} Pitching a model or a price before
  extracting the five numbers commits us to an answer we cannot yet
  defend.
\item
  \textbf{Accepting aspirations as SLOs.} ``Make it fast'' is a goal;
  ``p95 TTFT under 2 s'' is a contract. Write the contract.
\item
  \textbf{Letting the customer choose the technology.} ``We want a 70B
  model'' is a solution, not a requirement; the architect derives the
  model from the bounds (Ch5), not from the customer's guess.
\item
  \textbf{Ignoring the operational constraint.} ``It has to be on-prem /
  private / no cloud'' often decides more than any performance number
  --- surface it early.
\item
  \textbf{Not logging assumptions.} Every silence becomes a hidden
  assumption that surfaces as a dispute later. Write them down as
  {[}VERIFY{]} items.
\end{itemize}

\section{Architecture Consequence}\label{architecture-consequence}

The requirements conversation is where the book's entire loop is
\emph{fed.} The five bounds it extracts are exactly the inputs Ch1-16
consume: traffic feeds capacity (Ch17), tokens feed KV/memory (Ch7),
quality feeds model selection (Ch5), latency feeds serving (Ch11), cost
feeds TCO (Ch16), and the operational constraint picks the tier (Ch13).
A good requirements session is not a warm-up to architecture --- it
\emph{is} the first chapter of the architecture, and every downstream
decision traces back to the numbers it locked. {[}2°{]}

\begin{figure}
\centering
\pandocbounded{\includegraphics[keepaspectratio,alt={\hyperref[fig:23.1]{Fig 23.1} --- From vague ask to architectural bounds {[}ILLUSTRATIVE conceptual{]}},width=\maxwidth{\textwidth},keepaspectratio]{figures/fig-23-2301.pdf}}
\caption{\hyperref[fig:23.1]{Fig 23.1} --- From vague ask to architectural bounds
{[}ILLUSTRATIVE conceptual{]}}

\label{fig:23.1}\end{figure}

\emph{\hyperref[fig:23.1]{Fig 23.1} --- The framing funnel: broad ask → five Socratic
questions (problem, success, data, constraints, scope) → confirmed
bounds. This is the }dialogue* that surfaces requirements; the
\emph{six} architectural surfaces in \hyperref[tab:23.1]{Table 23-1} are what each answer
reveals (traffic, token profile, SLO, quality, ops/tier, TCO).
{[}ILLUSTRATIVE conceptual{]}*

\section{What We Still Don't Know}\label{what-we-still-dont-know}

\begin{itemize}
\tightlist
\item
  \textbf{How accurate the customer's own-told numbers are} (``40 rps
  peak'') is {[}VERIFY{]}; peak estimates are notoriously optimistic and
  should be triangulated with real logs where possible.
\item
  \textbf{Whether the quality bar once defined stays stable} is
  {[}HYPOTHESIS{]}; as users adopt the tool, their expectations and the
  data they query both shift.
\item
  \textbf{The right number of questions to lock quality} is
  {[}VERIFY{]}; 50 worked here as a representative proxy but a larger,
  task-distributed set is stronger.
\end{itemize}

\section{End-of-Chapter Mini-Case}\label{end-of-chapter-mini-case}

A CTO emails the architect: ``We heard LLMs are great --- let's put one
on our customer portal to answer support tickets. Shipping it next
month.'' Rather than confirm the launch date, the architect books a
30-minute session and runs the funnel. It emerges the portal gets
\textasciitilde5 rps average but 200-rps bursts during outage events
(traffic), tickets are short (\textasciitilde500 tokens, not
retrieval-heavy), quality means ``resolve the ticket, not just answer''
(a capability the cheaper model screens out), the SLO is 3 s because
users are already frustrated (latency), and --- critically --- the
portal sits on regulated payment data so nothing leaves the premises
(ops). The mirrored restatement turns ``let's ship AI next month'' into
``a 7B-class model on one on-prem GPU, resolving baseline tickets under
3 s, with a 200-rps burst buffer, quality-gated on real resolved tickets
--- and a two-week pilot on 5\% of tickets to confirm the referral
quality before full rollout.'' The customer's vague excitement became a
defensible, bounded, shippable scope --- and the architect is already
architecting, not guessing.
